%! The Idea of Elimination — Unit 2, Lesson 2.1 — Warm-Up (Answer Key)
\documentclass[10pt]{article}
\usepackage{linalg-article}
\usepackage{linalg-key}   % pulls in -boxes; do NOT also load -boxes
\newcommand{\TallMath}[1]{$\displaystyle #1\rule[-1.4em]{0pt}{3.2em}$}
\newcommand{\xx}{\mathbf{x}}
\newcommand{\bb}{\mathbf{b}}

\begin{document}

\pageheader{Unit 2, Lesson 2.1}{Warm-Up --- Answer Key}
\namedateperiod

\vspace{0.12in}

\begin{notesbox}{Getting Ready: Skills We'll Use Today}
\small These three moves are the pieces elimination is built from. Answer quickly.

\vspace{4pt}
\textbf{1. Subtract to remove a variable.} Both equations already have a $2y$. Subtract the first
from the second to remove $y$, then finish.
\[
  x + 2y = 7, \qquad 3x + 2y = 13.
\]
\hspace{2em} $x = {}$\ans{$3$} \qquad $y = {}$\ans{$2$}
\quad{\footnotesize(subtract: $2x=6$, so $x=3$, then $3+2y=7$, $y=2$)}

\vspace{8pt}
\textbf{2. Scale an equation.} Multiplying \emph{both sides} of an equation by the same number
keeps it true. Double every term of $2x + 3y = 13$:
\[
  2\,(2x + 3y) = 2\,(13) \quad\Longrightarrow\quad
  {\color{keyred}\mathbf{4}}\,x + {\color{keyred}\mathbf{6}}\,y = {\color{keyred}\mathbf{26}}.
\]

\vspace{4pt}
\textbf{3. Rows of $A\xx=\bb$.} Write $A\xx=\bb$ as two equations (one per row), with
$A=\begin{bmatrix} 2 & 3 \\ 4 & 7 \end{bmatrix}$, $\xx=\begin{bmatrix} x \\ y \end{bmatrix}$, and
$\bb=\begin{bmatrix} 13 \\ 27 \end{bmatrix}$.
\ansline{$2x + 3y = 13$;\qquad $4x + 7y = 27$}
\end{notesbox}

\begin{teachernote}
Item~1 is exactly the Lesson~2.0 move --- it works only \emph{because} both equations carry a
$2y$. Item~2 is the new ingredient: scaling a row so coefficients line up. Item~3 is the very
system solved in the notes; point out in the debrief that its $x$-coefficients ($2$ and $4$) do
\emph{not} match, so item~1's plain subtraction won't work --- that is why we need a multiplier.
\end{teachernote}

\end{document}
